% =====================================================================
%  SCL 模板（LaTeX + minted）
%  来源：Nemesis 队 Standard Code Library 的 main.tex
%  本文件保留了原始 main.tex 的全部“格式设置”（宏包、字体、minted 样式、
%  双栏、目录、页眉页脚、标题样式等），仅删除了所有原始内容。
%  你只需要在 \begin{document} ... \end{document} 之间按注释添加自己的内容。
%  编译方式：xelatex -shell-escape main.tex （详见 README.md / build.sh）
% =====================================================================

\documentclass[titlepage, a4paper]{article}
\usepackage{graphicx, amssymb, amsmath, textcomp, booktabs}
\usepackage[libertine,vvarbb]{newtxmath}
\usepackage[scr=rsfso]{mathalfa}
\usepackage[T1]{fontenc} % best for Western European languages
\usepackage{minted}
\usepackage{listings, color, setspace, titlesec, fancyhdr, mdframed, multicol}
\usepackage{ucharclasses}
\usepackage{xunicode, xltxtra}
\usepackage[inner=1.25cm, outer=0.8cm, top=1.7cm, bottom=0.0cm]{geometry}
\usepackage{pdfpages}
\usepackage{tocloft}
\usepackage{nameref}
\usepackage{verbatim}
\usepackage{relsize}
\usepackage{fontspec}
\usepackage[colorlinks, linkcolor = black]{hyperref}
\usepackage[table]{xcolor}
\usepackage{tabularx}
\usepackage{enumitem}
% configure fonts
% if not using CJK
% \newfontfamily\substitutefont{SimSun}[Scale=0.8,BoldFont=SimHei]
% \setTransitionsForChinese{\begingroup\substitutefont}{\endgroup}
\usepackage{xeCJK}
\setCJKmainfont{Noto Serif CJK SC}[Scale=0.8]
\setCJKmonofont{Noto Sans Mono CJK SC}[Scale=0.8]
\setCJKsansfont{Noto Sans CJK SC}[Scale=0.8]

\setmainfont{Linux Libertine O}[Scale=0.925]
\setmonofont{JetBrains Mono}[Scale=0.713]
%\setsansfont{Gill Sans Medium}

\XeTeXlinebreaklocale "zh"
\XeTeXlinebreakskip = 0pt plus 1pt

\setlength{\parindent}{0em}\setlength{\parskip}{1pt}
\setlength\itemsep{1pt}

\makeatletter
\renewcommand{\paragraph}{%
  \@startsection{paragraph}{4}%
  {\z@}{1pt \@plus 1pt \@minus 1pt}{-1em}%
  {\normalfont\normalsize\bfseries}%
}
\makeatother


%configure the top corners
\pagestyle{fancy}
\setlength{\headsep}{0.1cm}

\chead{SYNWYXCWL}
\rhead{Page \thepage}
\lhead{Shenzhen University}
 
%configure space between the two columns
\setlength{\columnsep}{13pt}

%configure minted to display codes
%\definecolor{Gray}{rgb}{0.9,0.9,0.9}

%remove leading numbers in table of contents
%\setcounter{secnumdepth}{0}	

%configure section style of table of content
\renewcommand\cftsecfont{\Large}

%configure section style
\titleformat{\section}
{\huge}			% The style of the section title
{\thesection.}				% a prefix
{4pt}						% How much space exists between the prefix and the title
{}					% How the section is represented
% \titleformat{\section}{\huge}{}{0pt}{}
\titlespacing{\section}{0pt}{0pt}{0pt}
\titlespacing{\subsection}{0pt}{0pt}{0pt}
\titlespacing{\subsubsection}{0pt}{0pt}{0pt}

%enable section to start new page automatically
%\let\stdsection\section
%\renewcommand\section{\penalty-100\vfilneg\stdsection}

%\renewcommand\theFancyVerbLine{\arabic{FancyVerbLine}}
\renewcommand{\theFancyVerbLine}{\sffamily \textcolor[rgb]{0.5,0.5,0.5}{\scriptsize {\arabic{FancyVerbLine}}}}

\setminted[cpp]{
	style=xcode,
	mathescape,
	linenos,
	autogobble,
	baselinestretch=0.8,
	tabsize=3,
	fontsize=\normalsize,
	%bgcolor=Gray,
	frame=single,
	framesep=1mm,
	framerule=0.3pt,
	numbersep=1mm,
	breaklines=true,
	breaksymbolsepleft=2pt,
	%breaksymbolleft=\raisebox{0.8ex}{ \small\reflectbox{\carriagereturn}}, %not moe!
	%breaksymbolright=\small\carriagereturn,
	breakbytoken=false,
	showtabs=true,
	tab={\relscale{0.6} $\big\vert \ \ \ $ \relscale{1}},
}
\setminted[java]{
	style=xcode,
	mathescape,
	linenos,
	autogobble,
	baselinestretch=0.8,
	tabsize=3,
	fontsize=\small,
	%bgcolor=Gray,
	frame=single,
	framesep=1mm,
	framerule=0.3pt,
	numbersep=1mm,
	breaklines=true,
	breaksymbolsepleft=2pt,
	%breaksymbolleft=\raisebox{0.8ex}{ \small\reflectbox{\carriagereturn}}, %not moe!
	%breaksymbolright=\small\carriagereturn,
	breakbytoken=false,
	showtabs=true,
	tab={\relscale{0.6} $\big\vert \ \ \ $ \relscale{1}},
}
\setminted[python]{
	style=xcode,
	mathescape,
	linenos,
	autogobble,
	baselinestretch=0.8,
	tabsize=3,
	fontsize=\small,
	%bgcolor=Gray,
	frame=single,
	framesep=1mm,
	framerule=0.3pt,
	numbersep=1mm,
	breaklines=true,
	breaksymbolsepleft=2pt,
	%breaksymbolleft=\raisebox{0.8ex}{ \small\reflectbox{\carriagereturn}}, %not moe!
	%breaksymbolright=\small\carriagereturn,
	breakbytoken=false,
	showtabs=true,
	tab={\relscale{0.6} $\big\vert \ \ \ $ \relscale{1}},
}

\setminted[vim]{
	style=xcode,
	mathescape,
	linenos,
	autogobble,
	baselinestretch=0.8,
	tabsize=2,
	fontsize=\normalsize,
	%bgcolor=Gray,
	frame=single,
	framesep=1mm,
	framerule=0.3pt,
	numbersep=1mm,
	breaksymbolsepleft=2pt,
	%breaksymbolleft=\raisebox{0.8ex}{ \small\reflectbox{\carriagereturn}}, %not moe!
	%breaksymbolright=\small\carriagereturn,
	breakbytoken=false,
}

\setminted[sh]{
	style=xcode,
	mathescape,
	linenos,
	autogobble,
	baselinestretch=0.8,
	tabsize=2,
	fontsize=\normalsize,
	%bgcolor=Gray,
	frame=single,
	framesep=1mm,
	framerule=0.3pt,
	numbersep=1mm,
	breaklines=true,
	breaksymbolsepleft=2pt,
	%breaksymbolleft=\raisebox{0.8ex}{ \small\reflectbox{\carriagereturn}}, %not moe!
	%breaksymbolright=\small\carriagereturn,
	breakbytoken=false,
}


%THE SCL BEGINS
\begin{document}
	\begin{multicols}{2}
		\setcounter{tocdepth}{3}
		\begingroup
		\let\cleardoublepage\relax
		\let\clearpage\relax
		\begin{small}
		\begin{spacing}{0.70}
		\tableofcontents
		\end{spacing}
		\end{small}
		\newpage
		\begin{spacing}{0.6}

			% ============================================================
			% 在这里添加你的内容。三种内容类型：
			%
			% 1) 代码：用 \inputminted 引入：
			%      \section{分类名}
			%        \subsection{算法名}
			%          \inputminted{cpp}{src/src/example/example.cpp}
			%    · 可选参数 highlightlines={7} 高亮某几行
			%
			% 2) 公式/笔记：写成独立 .tex 片段，用 \input 引入：
			%      \input{src/src/example/example.tex}
			%
			% 3) 正文内联公式：$...$ 或 \[ ... \]
			%
			% 4) 代码注释里的公式（已开启 mathescape）：
			%      在 .cpp 注释里直接写 // $O(n^2)$ 即可被渲染成数学式
			% ============================================================

			\section{Tips}
				多使用\texttt{\#define debug(...) fprintf(stderr, \_\_VA\_ARGS\_\_)}

			\section{Data Structure}
				\subsection{树状数组}
					\inputminted{cpp}{./BIT.cpp}
				\subsection{珂朵莉树}
					\inputminted{cpp}{./odt.cpp}
			
			\section{Math}
				\subsection{实数多项式}
					\inputminted{cpp}{./poly.cpp}
				\subsection{模多项式(NTT)}
					\inputminted{cpp}{./polyMod.cpp}
				\subsection{插值}
					\begin{multicols}{2}
    Lagrange
    
    $
        f(x) = \sum_{i=1}^{n}y_i\prod_{j\neq i}\frac{x - x_j}{x_i - x_j}
    $

    暴力处理 $O(n^2)$, 可用多项式优化

    $x_i$ 取连续整数时可优化至 $O(n)$

    \newcolumn
    Newton

\end{multicols}

					\inputminted{cpp}{./interp_lagrange.cpp}
			
			\section{String}
				\subsection{SAM 后缀自动机}
					\inputminted{cpp}{./sam.cpp}
			
			\section{Geometry}
				\subsection{整点凸包}
					\inputminted{cpp}{./geometry.cpp}
			

		\end{spacing}
		\endgroup
	\end{multicols}
\end{document}
%THE SCL ENDS
